\documentclass[11pt]{article}
\usepackage{preamble}
\usepackage{multicol}
\setlength{\parskip}{4pt}

\begin{document}

\daybanner{AP Statistics \textperiodcentered\ Unit 1 \textperiodcentered\ Topics 1.4--1.5 (CED 1.4--1.5) \textperiodcentered\ E: 2026-09-16 \textperiodcentered\ TEACHER KEY}{One Commute, Two Questions}

\sectionbanner{CED TETHER}
\begin{multicols}{2}
{\footnotesize\setlength{\parskip}{0pt}
\begin{itemize}[leftmargin=*,itemsep=0pt,topsep=0pt]
\item Enduring Understanding: UNC-1 Graphical representations and statistics allow us to identify and represent key features of data.
\item Skills:
\item \textbf{Skill 2.B} --- Construct numerical or graphical representations of distributions.
\item \textbf{Skill 2.A} --- Describe data presented numerically or graphically.
\item \textbf{Skill 2.D} --- Compare distributions or relative positions of points within a distribution.
\item Learning Objectives and Essential Knowledge:
\item LO UNC-1.C Represent categorical data graphically. [Skill 2.B]
\item EK UNC-1.C.1 Bar charts (or bar graphs) are used to display frequencies (counts) or relative frequencies (proportions) for categorical data.
\item EK UNC-1.C.2 The height or length of each bar in a bar graph corresponds to either the number or proportion of observations falling within each category.
\item EK UNC-1.C.3 There are many additional ways to represent frequencies (counts) or relative frequencies (proportions) for categorical data.
\item LO UNC-1.D Describe categorical data represented graphically. [Skill 2.A]
\item EK UNC-1.D.1 Graphical representations of a categorical variable reveal information that can be used to justify claims about the data in context.
\item LO UNC-1.E Compare multiple sets of categorical data. [Skill 2.D]
\item EK UNC-1.E.1 Frequency tables, bar graphs, or other representations can be used to compare two or more data sets in terms of the same categorical variable.
\item \textbf{EU UNC-1} --- Graphical representations and statistics allow us to identify and represent key features of data.
\item UNC-1.F Classify types of quantitative variables. [Skill 2.A]
\item UNC-1.F.1 A discrete variable can take on a countable number of values. The number of values may be finite or countably infinite, as with the counting numbers.  UNC-1.F.2 A continuous variable can take on infinitely many values, but those values cannot be counted. No matter how small the interval between two values of a continuous variable, it is always possible to determine another value between them.
\item UNC-1.G Represent quantitative data graphically. [Skill 2.B]
\item UNC-1.G.1 In a histogram, the height of each bar shows the number or proportion of observations that fall within the interval corresponding to that bar. Altering the interval widths can change the appearance of the histogram.  UNC-1.G.2 In a stem and leaf plot, each data value is split into a "stem" (the first digit or digits) and a "leaf" (usually the last digit).  UNC-1.G.3 A dotplot represents each observation by a dot, with the position on the horizontal axis corresponding to the data value of that observation, with nearly identical values stacked on top of each other.  UNC-1.G.4 A cumulative graph represents the number or proportion of a data set less than or equal to a given number.  UNC-1.G.5 There are many additional ways to graphically represent distributions of quantitative data.
\end{itemize}}
\end{multicols}
\clearpage
\small
\textbf{Essential question:} Which questions can each display answer, and which connection is lost?\par
\textbf{DOK-3 skill:} 4.B \textperiodcentered\ \textbf{FRQ pattern:} {\small\texttt{adjudicate-\allowbreak{}competing-\allowbreak{}displays-\allowbreak{}and-\allowbreak{}bound-\allowbreak{}the-\allowbreak{}report}}\par
\textbf{DOK rationale:} Part (c) coordinates the two topics to judge competing claims, defend a corrected report, and explain what the rejected representation cannot establish.\par

\frameworkphaseheader{First take $\rightarrow$ both follow-alongs $\rightarrow$ finish}{1 $\rightarrow$ 3}{43.8}{%
\begin{itemize}[leftmargin=1.5em,itemsep=2pt,topsep=2pt]
  \item Read the competing reports without endorsing either. Start the first take at the bell.
  \item Play both topic videos in teaching order; pause for each follow-along.
  \item Allow ten minutes to finish (a)--(c). Collect one sheet; score (c) only, E/P/I by a human.
\end{itemize}
}{%
\begin{itemize}[leftmargin=1.5em,itemsep=2pt,topsep=2pt]
  \item Commit one sentence before watching.
  \item Complete both follow-alongs, then finish this single problem and turn it in.
\end{itemize}
}{%
\begin{itemize}[leftmargin=1.5em,itemsep=2pt,topsep=2pt]
  \item Which feature settles that claim?
  \item What information would your chosen display hide?
\end{itemize}
}{Circulate and ask for evidence linked to a claim. This is the human-graded DOK-3 channel; do not send it to an AI grader or auto-score it.}

\begin{callout}{\IconWarn\ WATCH FOR}{calloutred}
\begin{itemize}[leftmargin=1.5em,itemsep=2pt,topsep=2pt]
  \item Calling Bus the unique winner despite the tie.
  \item Treating a histogram interval as one exact travel time.
  \item Assigning long trips to bus riders without linked observations.
\end{itemize}
\end{callout}

\sectionbanner{SCORING PART (c) --- E / P / I}

\textbf{Expected elements}
\begin{itemize}[leftmargin=1.5em,itemsep=2pt,topsep=2pt]
  \item \textbf{e1} (required) --- Uses the tied counts to bound the most-common-mode claim.
  \item \textbf{e2} (required) --- Assigns categorical comparison to the bar chart and quantitative distribution to the histogram, with a numerical claim.
  \item \textbf{e3} (required) --- Explains that pooled displays lack individual mode--time links; identifies the needed paired information.
\end{itemize}

{\small\renewcommand{\arraystretch}{1.25}
\noindent\begin{tabular}{|p{0.5in}|p{5.9in}|}
\hline
\textbf{E} & All three evidence-to-reason links are correct in context. Accept an equivalent defensible report. \\ \hline
\textbf{P} & At least one correct evidence-to-reason link, but another is missing or incorrect. \\ \hline
\textbf{I} & No correct evidence-to-reason link; only a choice, copied numbers, or unsupported claims. \\ \hline
\end{tabular}}

\textbf{Common mistakes}
\begin{itemize}[leftmargin=1.5em,itemsep=2pt,topsep=2pt]
  \item Calling Bus the unique winner despite the tie.
  \item Treating a histogram interval as one exact travel time.
  \item Assigning long trips to bus riders without linked observations.
\end{itemize}

\clearpage
\normalsize
\focusbanner{TODAY'S PROBLEM --- annotated}

Suppose ten students report their usual travel mode and travel time. Both graphs summarize these same students; the individual mode--time links are not shown. Histogram intervals include the left endpoint only.\par\smallskip \textbf{Ari:} ``Bus is the most common mode, so the bus riders must be the students with long trips. '' \textbf{Dev:} ``The histogram answers every question about these commutes; we can drop the mode graph.''\par
\begin{center}
\begin{tikzpicture}
\begin{axis}[width=0.96\linewidth,height=1.35in,ybar,bar width=16pt,ymin=0,ymax=5,
 symbolic x coords={Walk,Bus,Car},xtick=data,ytick={0,2,4},ylabel={Students},
 title={Travel mode},tick label style={font=\scriptsize},label style={font=\scriptsize},title style={font=\small}]
\addplot[fill=calloutblue,draw=dayblue] coordinates {(Walk,4) (Bus,4) (Car,2)};
\end{axis}\end{tikzpicture}\par\smallskip
\begin{tikzpicture}
\begin{axis}[width=0.96\linewidth,height=1.35in,ybar interval,x tick label as interval=false,bar shift=0pt,xmin=0,xmax=30,
 ymin=0,ymax=6,xtick={0,10,20,30},ytick={0,2,4,6},xlabel={Travel time (minutes)},ylabel={Students},
 tick label style={font=\scriptsize},label style={font=\scriptsize}]
\addplot[fill=calloutblue,draw=dayblue] coordinates {(0,5) (10,2) (20,3) (30,0)};
\end{axis}\end{tikzpicture}
\end{center}
\begin{firsttakebox}{FIRST TAKE (5 min)}
Which report would you send to the principal, or would you revise both? Name one visible feature.\par
\answer{Ungraded commitment; accept any evidence-based first impression.}
\end{firsttakebox}

\textit{Rules callout ``TWO TOPICS, ONE DATA STORY'' is printed on the student sheet.}\par\medskip

\dokbadge{1}~\textbf{(a)} Name the most common travel mode or modes and their counts.\par
\answer{Walk and Bus tie at four students each; neither is the unique most common mode.}

\dokbadge{2}~\textbf{(b)} How many students have trips of at least 10 minutes? Explain how the histogram bars answer this and why they do not give each exact travel time.\par
\answer{Two plus three gives five students with trips of at least 10 minutes. A histogram groups observations into intervals; the two students in $[10,20)$ need not have the same travel time.}

\dokbadge{3}~\textbf{(c)} $\star$ In three to five sentences, judge both reports. Choose a display for comparing travel modes and a display for describing travel times. Give one corrected numerical claim and explain what extra information is needed to connect long trips with bus travel.\par
\answer{Revise both. The bar chart compares modes and shows Walk and Bus tied at four, so Ari cannot claim a unique bus lead or identify the long trips as bus trips. The histogram describes the quantitative times and shows five of ten trips last at least 10 minutes, but it does not preserve mode labels. Keep both for the two questions; matching each student\textquotesingle s mode with their time is needed to compare long trips across modes.}

\begin{summaryexitbox}{EXIT}
One claim I would revise after both videos: \hrulefill
\end{summaryexitbox}

\end{document}
